\scanpage{13}
On utilise un corollaire nouveau \unsure{?|:}
\margin{\unsure{!!}}

\margin{$\overline{\text{\uline{XI}}}$ \deleted{\unsure{XII|4}}}
\uline{Corollaire} 4.\deleted{\unsure{8}}\unsure{7|8}. \deleted{\uline{\unsure{base locale}}}

$G$ lisse affine sur $S$ local de pt fermé $s$. $\deleted{\unsure{H}\ \unsure{(H}}$ $H(n)$ famille de ss-gpes de $G$, de t. m., $H(n)$, $={}_nH(m)$, on suppose qu'il existe un ss-gpe de t. m. $H_s$ de $G_s$ \unsure{tels|tel} que $H(n)_s={}_nH_s$ pour tt $n$. Alors $\exists$ ss-gpe $H$ \inserted{de t. m.} unique de $G$ tel que $ {}_nH\unsure{=|\mathrel{:=}} \unsure{H_n|H(n)}$, $\forall n$.

\uline{Remarque} \deleted{Il suffit} Enoncé analogue en se bornant, aux $n$ dont l'un des diviseurs premiers est $\in P\ne\varnothing$ ($P$ ens de \unsure{nbres|nbs} premiers) pourvu que la torsion du type $H_s$ soit première à $P$. On utilise pour ceci la variante de la démonstration précédente, avec $M$ \unsure{courant|parcourant|remplacé} un \deleted{\unsure{sys}} ss-groupe de t. m. de $=$ type que $H_s$ [et en \deleted{\unsure{se bornant}} utilisant le th. de densité dans la dém. sous la forme \ldots]
% archive sheet number: 12
